% =====================
\section{検証と妥当性確認：V\&V計画}

ハッカソン2テキスト§1.2.5に基づき，MOE → MOP → TPMの順で評価指標を定める．
実運用では順序が逆転し，実装中にTPMを継続監視し，V\&VによってMOPを確認し，
最終的にMOEで受入判定を行う．

\subsection{プロジェクトの成果物に対する有効指標MOE}

MOE（Measure of Effectiveness：有効指標）は，ステークホルダ（審査員・聴衆）が
成果物の合格・不合格を判定するための指標であり，受入テストの合否基準に相当する．

\begin{table}[H]
  \centering
  \begin{tabularx}{\linewidth}{clXl}
    \toprule
    No. & MOE & 内容 & 確認方法 \\
    \midrule
    1 & 演奏完走 &
      「かえるのうた」5パート輪唱（テンポ変化・フェルマータ含む）を
      最初から最後まで演奏できること &
      実演で確認 \\
    2 & 指揮追従 &
      指揮者人形のテンポ変化・フェルマータに対して
      演奏が連動していることが視聴覚的に確認できること &
      審査員観察評価 \\
    3 & 楽器音認識 &
      各パートが意図した楽器（ピアノ・鉄琴・トランペット・
      バス/スネア・ハイハット）の音として聴こえること &
      審査員主観評価 \\
    \bottomrule
  \end{tabularx}
  \caption{MOE一覧}
  \label{tab:moe}
\end{table}

\subsection{有効指標MOEに対応する性能指標MOP}

MOP（Measure of Performance：性能指標）はMOEを達成するために必要な
システム性能を定量的に示す指標であり，V字モデルの設計段階で定義する．
FBSの各機能に対応する性能要素に対して設定した．

\begin{table}[H]
  \centering
  \begin{tabularx}{\linewidth}{clXl}
    \toprule
    No. & 対応MOE & MOP & 対応テスト \\
    \midrule
    1 & MOE-1,2 & フォトトランジスタによる拍検出成功率 & 単体テスト \\
    2 & MOE-1,2 & 楽器間の発音タイミング誤差 & 結合テスト \\
    3 & MOE-1,2 & テンポ変化後のEMA収束拍数 & システムテスト \\
    4 & MOE-1,2 & 腕停止からフェルマータ状態移行までの拍数 & システムテスト \\
    5 & MOE-3   & 各パートの発音周波数精度 & 単体テスト \\
    \bottomrule
  \end{tabularx}
  \caption{MOEとMOPの対応}
  \label{tab:moe-mop}
\end{table}

\subsection{システムテストで評価する性能指標MOP}

\begin{table}[H]
  \centering
  \begin{tabularx}{\linewidth}{llXl}
    \toprule
    No. & 対応FBS機能 & MOP（目標値） & 対応テスト \\
    \midrule
    SYS-1 & テンポ推定（F2a3） &
      テンポ変化後の収束：$\leq 3$拍 & テストS1 \\
    SYS-2 & フェルマータ検出（F2a4） &
      腕停止からフェルマータ移行：$\leq 1$拍 & テストS2 \\
    SYS-3 & 楽譜進行制御（F2c2） &
      全区間通し演奏の完走率：100\% & テストS3 \\
    \bottomrule
  \end{tabularx}
  \caption{システムテストで評価するMOP}
  \label{tab:mop-sys}
\end{table}

\subsection{結合テストで評価する性能指標MOP}

\begin{table}[H]
  \centering
  \begin{tabularx}{\linewidth}{llXl}
    \toprule
    No. & 対応FBS機能 & MOP（目標値） & 対応テスト \\
    \midrule
    INT-1 & テンポ追従演奏（F2c3） &
      楽器間発音タイミング誤差：$\leq \pm 20\,\mathrm{ms}$ & テストC2 \\
    INT-2 & テンポ追従演奏（F2c3） &
      Arduino→Processing転送遅延の計測と補正値確定 & テストC1 \\
    INT-3 & 輪唱オフセット制御（F2c1） &
      輪唱オフセット誤差：$\leq \pm 1$拍 & テストC3 \\
    \bottomrule
  \end{tabularx}
  \caption{結合テストで評価するMOP}
  \label{tab:mop-int}
\end{table}

\noindent
INT-1の目標値$\pm 20\,\mathrm{ms}$は，人間が音楽的なタイミングのズレを知覚できる
閾値（先行研究では概ね$20\sim30\,\mathrm{ms}$）を参考として設定した．
確定値は結合テスト（テストC2）の実測結果をもって設定する．

\subsection{単体テストで評価する性能指標MOP}

\begin{table}[H]
  \centering
  \begin{tabularx}{\linewidth}{llXl}
    \toprule
    No. & 対応PBS成果物 & MOP（目標値） & 対応テスト \\
    \midrule
    UNIT-1 & フォトTr配線（PB1b） &
      拍検出成功率：$\geq 95\,\%$ & テストP2/P3 \\
    UNIT-2 & 赤外線センサー配線（PB1a） &
      距離計算誤差：$\leq \pm X\,\mathrm{cm}$（実測で確定） & テストP1 \\
    UNIT-3 & 状態機械SW（PB2b） &
      EMA通常演奏時（$\alpha=0.2$）の収束：$\leq 13$拍 & テストU1 \\
    UNIT-4 & 演奏制御SW（PC2a） &
      NOTE\_ON送信タイミング誤差：$\leq \pm X\,\mathrm{ms}$（実測で確定） & テストU3 \\
    UNIT-5 & 演奏制御SW（PC2a） &
      発音周波数精度：目標周波数の$\pm 1\,\mathrm{Hz}$以内 & テストU4 \\
    \bottomrule
  \end{tabularx}
  \caption{単体テストで評価するMOP}
  \label{tab:mop-unit}
\end{table}

\subsection{プロジェクト中に監視する技術性能指標TPM}

TPM（Technical Performance Measure：技術性能指標）は，開発・実装中に
継続的に測定・報告する定量的指標であり，技術的リスクの早期発見を目的とする．
MOPのうち特にシステム完成時の性能に大きく影響する項目を選定した．

\begin{table}[H]
  \centering
  \begin{tabularx}{\linewidth}{llXl}
    \toprule
    No. & 対応PBS要素 & TPM（監視値・許容範囲） & 確定方法 \\
    \midrule
    TPM-1 & 赤外線センサー（PB1a） &
      距離計算誤差：$\leq \pm X\,\mathrm{cm}$ & 単体テストP1 \\
    TPM-2 & フォトTr（PB1b） &
      LED検出閾値電圧（実機測定値） & 単体テストP2/P3 \\
    TPM-3 & 演奏制御SW（PC2a） &
      Arduino→Processing転送遅延：$\leq X\,\mathrm{ms}$ & 結合テストC1 \\
    TPM-4 & 状態機械SW（PB2b） &
      EMA（$\alpha=0.2$）テンポ収束拍数：$\leq 13$拍 & 単体テストU1 \\
    TPM-5 & サーボ機構（PA1a） &
      指揮者人形が追従できる最大BPM：$\geq 104\,\mathrm{bpm}$ & 単体テストP4/P5 \\
    \bottomrule
  \end{tabularx}
  \caption{TPM一覧}
  \label{tab:tpm}
\end{table}

\noindent
$X$と記載した値は実機測定・試作時に確定する．確定後は本表を更新すること．

\subsection{プロジェクト中に監視する指標TPM}

各TPMの監視頻度・方法・担当を表\ref{tab:tpm-monitor}に示す．

\begin{table}[H]
  \centering
  \begin{tabularx}{\linewidth}{llXl}
    \toprule
    TPM & 監視頻度 & 監視方法 & 担当 \\
    \midrule
    TPM-1 & 試作時・単体テスト時 &
      センサー実測値と近似式計算値の誤差を記録 & B \\
    TPM-2 & 試作時・単体テスト時 &
      シリアルモニタまたはオシロスコープで閾値電圧を確認 & A/B \\
    TPM-3 & 結合テスト時 &
      送信タイムスタンプと受信タイムスタンプの差分を計測 & B/C \\
    TPM-4 & 単体テスト時 &
      シミュレーションデータを用いてEMAの収束拍数を計測 & B \\
    TPM-5 & 試作時・単体テスト時 &
      指定BPMで人形を動作させ，遅延・脱調の有無を確認 & A \\
    \bottomrule
  \end{tabularx}
  \caption{TPM監視体制}
  \label{tab:tpm-monitor}
\end{table}
